\section{Limitations and Discussion}\label{sec:discussion}
\textbf{Guiding model selection for IoT systems:} We believe that our work provides complementary evidence for selecting anomaly detectors, because it separates the effect of the framework from the effect of the evaluation protocol. Our results show that practitioners should fix and audit the protocol before they compare frameworks. The protocol changes the macro-F1 by up to 0.13, whereas the largest difference between the frameworks on MEMS is 0.013. Framework rankings obtained under a random split of time-ordered sensor data should therefore not guide a deployment decision.

\textbf{Evaluation protocol:} A held-out block corresponds to a recording condition that the model has not observed, so that the block protocol measures generalization to new conditions. A random protocol places neighboring samples of one recording in both partitions, and test samples then have close neighbors in the training data. The scores of different methods rose and fell together across seeds, which indicates that the choice of held-out blocks decided most of the variation.

\textbf{AutoML versus default baselines:} A default random forest stayed within 0.035 macro-F1 of both frameworks in every setup and matched or exceeded them on deduplicated device 9. This agrees with the extensive benchmark of~\cite{gijsbers2022}, where average ranks of frameworks rarely differed significantly. Two explanations are plausible for this observation, and we did not test either of them. The feature sets contain only 3 and 20 features, so a tree ensemble can extract most of the available signal. In addition, a budget of 1800\,s may be short for a search over the full set of algorithm families.

The value of AutoML lies in its protection against a poor default. The default MLP was worse than both frameworks by 0.09 to 0.11 on deduplicated device 9, and the default decision tree was worse by 0.04 to 0.10 on MEMS. The frameworks avoided both outcomes without any manual tuning of the models. For a practitioner, this protection matters more than a small gain over the best default.

\textbf{Class balancing:} Under the block split, both frameworks detect the failure class least reliably, with a recall of 0.36 to 0.40, which is the least favorable outcome for predictive maintenance. Oversampling did not help, and it lowered the macro-F1 of both frameworks under the block split. A likely reason is that the oversampled minority samples originate from few recording blocks. Repeating them makes the model fit these blocks more closely, and it does not teach the model about unseen conditions. We did not test this explanation in the present study. Oversampling raised the recall of the failure class and lowered the recall of the normal class, so it traded false alarms for detected failures, and an application may accept this trade deliberately. A maintenance team may, for example, prefer to find more failures at the price of more false alarms. The built-in balancing of H2O AutoML collapsed the recall of the normal class in one seed, which shows that such an option needs repeated runs before it is trusted.

\textbf{Data quality:} On device 9, the ceiling of 0.718 and the near-identical scores of the best methods indicate that the representation limits the result. More search cannot recover information that the features do not contain. A plausible cause is that the 20-feature subset can map distinct traffic states, such as idle periods of different devices, to the same vector. We did not test this cause, because the full feature set is unavailable. Accuracy was misleading in the deduplicated variant (0.98 against a macro-F1 of 0.79), as in our earlier intrusion detection study~\cite{velnambi2026}. The order of the three best methods depended on about 34 test samples, and studies should therefore report the number of test samples per class next to every macro-averaged metric.

\textbf{Cost and deployment:} A random forest with default settings matched the frameworks on device 9. Its batch prediction was about 110 times faster than that of Auto-sklearn and 23 times faster than that of H2O AutoML. Table~\ref{tab:models} explains the size of the models of both frameworks. Random forests carry most of the ensemble weight of Auto-sklearn on device 9, and $k$-nearest-neighbor models, which store their training data, carry 45\% of the weight on MEMS with oversampling. The model size did not follow the number of ensemble members (correlation 0.29) but their type.

For IoT gateways, a model of 4.5\,GB is often unusable, and three options exist for deploying such a model on a gateway. Practitioners can restrict the search space and exclude forests with many trees and neighbor models. They can compress the ensemble after the search has finished. Alternatively, they can use the AutoML result only to select a model family and deploy a smaller model of that family.

\textbf{Feature reliance:} Two frameworks with different search strategies selected models that rely on the same few features. These features belong to the family that was found for hand-built models~\cite{gummadi2024xaiiot}, which suggests that the explanatory structure originates from the data and not from the learner. The feature \texttt{HH\_L1\_weight} describes the recent traffic volume of a host, so it likely separates floods from quiet traffic. We did not test this reading of the feature. Practitioners can trust the leading features across frameworks, but they should not interpret the order of features with near-zero importance.

\textbf{Scope of the datasets:} We evaluated one sensor dataset and device 9 of the traffic dataset in its released 20-feature form. The conclusions may not transfer to other devices or to the original 115 features. We did not map class identifiers to attack names, because the released files contain only integer identifiers.

\textbf{Inferred block structure:} The label blocks of MEMS stand in for sessions, because true session or speed identifiers are not available. The block protocol may therefore be easier or harder than a split by real sessions.

\textbf{Search budget and framework selection:} All runs used a budget of 1800\,s, and rankings could change with smaller or larger budgets. We compared two frameworks, and further frameworks such as FLAML~\cite{wang2019flaml} are natural candidates for future work. The pipeline accepts them without change, because it only requires a function that returns a fitted model for given training data, budget, threads, and seed.

\textbf{Statistical power:} Five seeds give coarse estimates of variance, and no paired rank test can reach $p<0.05$ with five seeds. Searches that stop on a wall-clock limit are also not exactly reproducible, even with fixed seeds.

\textbf{Different validation strategies:} H2O AutoML validates by cross-validation, and Auto-sklearn validates by holdout. The numbers of evaluated models are therefore not comparable, and Auto-sklearn exceeded the budget by up to 249\,s.

\textbf{Cost measurement:} The peak memory sums the resident set sizes of all processes and counts shared pages once per process. The size of the Auto-sklearn model is the size of the pickled classifier, which overstates its inference footprint. The latency of H2O AutoML includes the conversion of data into its frame format, and we do not report processor utilization.

\textbf{Explanation validity:} We computed permutation importance on 3,000 test rows of one device. Classes 1, 3, and 4 (about 26, 3, and 5 of roughly 35,000 test rows) are almost absent from this subset, so the importances describe decisions on the separable classes. Correlated window statistics can also distort permutation importance~\cite{strobl2008}. MEMS and the raw variant of device 9 were not analyzed with this method.

\textbf{Ambiguity ceiling:} The ceiling is computed on the full raw file and assumes that every non-ambiguous vector can be learned.

\textbf{Recommendations for practitioners:} Our results suggest six rules for studies and deployments. Audit the data for duplicates and ambiguous inputs before models are compared. Hold out entire sessions or blocks when the data are time ordered. Report macro-averaged metrics and per-class recall together with the number of test samples per class. Include default baselines, such as a random forest, in every comparison. Repeat every run with multiple seeds and report paired differences. Measure the size and latency of the model and compare them with the limits of the target device.
